\documentclass[12pt,a4paper]{article}

\usepackage[utf8]{inputenc}
\usepackage[T1]{fontenc}
\usepackage{amsmath, amssymb, amsthm}
\usepackage{geometry}
\usepackage{booktabs}
\usepackage[dvipsnames]{xcolor}
\usepackage{hyperref}
\usepackage{parskip}

\geometry{margin=2.5cm}

\hypersetup{colorlinks=true, linkcolor=blue!60!black, citecolor=blue!60!black, urlcolor=blue!60!black}

\newtheorem{hypothesis}{Hipotesis}
\newtheorem{remark}{Observacion}

\title{\textbf{Relacion entre el Termino Fuente $b(x)$\\
y el Operador Diferencial $\mathcal{L}$}\\[0.5em]
\large Teoria, Hipotesis y Literatura\\[0.3em]
\normalsize Notas para un Curso de Maestria en Boundary Element Method}
\author{W.F. Florez\\
\small Wessex Institute of Technology}
\date{}

\begin{document}

\maketitle
\tableofcontents
\newpage

% ─────────────────────────────────────────────────────────────
\section{El Problema Central}
% ─────────────────────────────────────────────────────────────

Dentro del marco de conversion de integrales de dominio a integrales de frontera
mediante el metodo de soluciones particulares, la pregunta fundamental es:

\medskip
\begin{center}
\textit{Dado $b(x)$, ?`cual es el operador diferencial $\mathcal{L}$ optimo tal
que la conversion sea eficiente, precisa y bien condicionada?}
\end{center}
\medskip

Recordemos que la conversion descansa sobre encontrar $\hat{\phi}_j$ tal que
$\mathcal{L}\hat{\phi}_j = f_j(r)$, y que la identidad de Green para $\mathcal{L}$
transforma:

\begin{equation}
    \int_\Omega b(x)\, d\Omega \approx
    \sum_j \alpha_j \int_\Omega \mathcal{L}\hat{\phi}_j\, d\Omega
    = \sum_j \alpha_j \int_\Gamma \mathcal{B}[\hat{\phi}_j]\, d\Gamma
    \label{eq:conversion}
\end{equation}

donde $\mathcal{B}$ es el operador de frontera emergente de la identidad de Green
asociada a $\mathcal{L}$. La pregunta sobre la eleccion optima de $\mathcal{L}$
involucra dos condiciones que no son independientes: que $b$ sea bien aproximable
por la base $\{f_j\} = \{\mathcal{L}\hat{\phi}_j\}$, y que la identidad de Green
para $\mathcal{L}$ convierta la integral de dominio en el menor numero posible de
terminos de frontera.

% ─────────────────────────────────────────────────────────────
\section{Las Dos Condiciones de Golberg et al.}
% ─────────────────────────────────────────────────────────────

La referencia mas precisa sobre los criterios de eleccion de las funciones base
es \cite{golberg1998comments}, donde se establecen dos condiciones necesarias que
la base $\{f_j\}$ debe satisfacer:

\begin{enumerate}
    \item \textbf{Aproximacion precisa:} $\{f_j\}$ debe proporcionar una
    aproximacion suficientemente precisa de $b(x)$ en los nodos de colocacion.
    Esto es un requisito de densidad de la base en un espacio funcional apropiado.

    \item \textbf{Integrabilidad analitica:} debe existir $\hat{\phi}_j$ en forma
    cerrada tal que $\mathcal{L}\hat{\phi}_j = f_j$. Sin esto, la conversion a
    integrales de frontera no es computable eficientemente.
\end{enumerate}

\begin{remark}
Estas dos condiciones son necesarias pero no suficientes para optimalidad.
No dicen nada sobre como elegir $\mathcal{L}$ dado $b$, sino que caracterizan
que pares $(\mathcal{L}, \{f_j\})$ son admisibles una vez fijado el operador.
\end{remark}

Adicionalmente, \cite{golberg1999rbf} senala una observacion critica:

\begin{quote}
\textit{Existe una interaccion sutil entre el error de aproximacion RBF y el
error BEM --- estos errores no son simplemente aditivos. Los estudios
experimentales que discuten solo errores RBF pueden llegar a conclusiones
erroneas sobre el error total en el DRM.}
\end{quote}

Este resultado subraya que la eleccion de $\mathcal{L}$ afecta no solo la
aproximacion de $b$ sino tambien la propagacion del error en toda la cadena
numerica.

% ─────────────────────────────────────────────────────────────
\section{El Metodo del Aniquilador: Marco Teorico}
% ─────────────────────────────────────────────────────────────

La conexion mas profunda entre $b(x)$ y $\mathcal{L}$ proviene de la teoria del
\textbf{metodo del aniquilador}, generalizado a EDPs por \cite{golberg1999annihilator}.

\subsection{Definicion de Aniquilador}

Sea $\mathcal{A}$ un operador diferencial lineal. Se dice que $\mathcal{A}$ es
\textbf{aniquilador} de $b(x)$ si:

\begin{equation}
    \mathcal{A}\, b(x) = 0
\end{equation}

En el caso de EDOs, si $b(x) = e^{\lambda x}$, entonces $\mathcal{A} = D - \lambda$
es el aniquilador, ya que $(D - \lambda)e^{\lambda x} = 0$. Para
$b(x) = \sin(kx)$, el aniquilador es $\mathcal{A} = D^2 + k^2$.

\subsection{Extension a EDPs: el Resultado de Golberg-Chen-Rashed}

\cite{golberg1999annihilator} generalizan este concepto a ecuaciones en derivadas
parciales. La idea central es:

\begin{equation}
    \mathcal{A}\, b = 0 \implies
    \text{usar } \mathcal{L} \text{ relacionado con } \mathcal{A}
    \text{ como operador de la conversion}
\end{equation}

La logica es que si $b$ pertenece al espacio nulo de $\mathcal{A}$, entonces
las funciones base $\{f_j\} = \{\mathcal{L}\hat{\phi}_j\}$ deben aproximar
funciones en ese mismo espacio nulo, lo cual ocurre naturalmente cuando
$\mathcal{L}$ esta relacionado con $\mathcal{A}$.

Especificamente, si $\mathcal{A} = \nabla^2 + k^2$ (operador de Helmholtz) y
$b$ es aproximadamente oscilatoria con numero de onda $k$, entonces usar
$\mathcal{L} = \nabla^2 + k^2$ como operador de la conversion es
\emph{compatible} con la estructura de $b$.

\subsection{Aplicacion a las RBFs Poliharmonicas}

\cite{golberg1999annihilator} aplican el metodo del aniquilador para encontrar
$\hat{\phi}$ de manera sistematica cuando $f_j$ son \emph{thin plate splines}
(TPS) y \emph{higher order splines}. El procedimiento es:

\begin{enumerate}
    \item Identificar el aniquilador $\mathcal{A}$ de $f_j(r)$.
    \item Aplicar $\mathcal{A}$ a $\mathcal{L}\hat{\phi} = f_j$, obteniendo
    $\mathcal{A}\mathcal{L}\hat{\phi} = \mathcal{A}f_j = 0$.
    \item La ecuacion homogenea resultante es solucionable por metodos clasicos,
    y comparando coeficientes se determina $\hat{\phi}$.
\end{enumerate}

Esto es una generalizacion del clasico metodo de coeficientes indeterminados
de ODEs, y proporciona una receta sistematica para obtener $\hat{\phi}$ sin
necesidad de adivinar su forma funcional.

% ─────────────────────────────────────────────────────────────
\section{Clasificacion de Operadores y sus Identidades de Green}
% ─────────────────────────────────────────────────────────────

La conversion \eqref{eq:conversion} requiere una identidad de Green para
$\mathcal{L}$. La complejidad del operador de frontera $\mathcal{B}$ resultante
depende del orden de $\mathcal{L}$. Se presenta una clasificacion de los casos
mas relevantes en la literatura.

\subsection{Operador de Laplace: $\mathcal{L} = \nabla^2$}

La identidad de Green de primer orden:

\begin{equation}
    \int_\Omega \nabla^2\hat{\phi}\, d\Omega
    = \int_\Gamma \frac{\partial\hat{\phi}}{\partial n}\, d\Gamma
\end{equation}

Un solo termino de frontera. Es la conversion mas limpia posible. Como muestra
\cite{cheng2000polyharmonic}, para el operador de Laplace las soluciones
particulares de funciones conicas y splines se obtienen por integracion
repetida en la coordenada radial --- lo que desarrollamos explicitamente al
derivar $\hat{\phi}$ para $f = 1 + r^3$:

\begin{equation}
    \frac{1}{r}\frac{d}{dr}\left(r\frac{d\hat{\phi}}{dr}\right) = f(r)
    \implies \hat{\phi}(r) = \frac{r^2}{4} + \frac{r^5}{25}
    \quad \text{para } f = 1 + r^3
\end{equation}

\textbf{Compatibilidad con $b$:} optimo cuando $b$ es una funcion suave sin
estructura oscilatoria ni de decaimiento pronunciado. En ausencia de
informacion sobre la naturaleza de $b$, $\nabla^2$ es la eleccion por
defecto.

\subsection{Operador de Helmholtz: $\mathcal{L} = \nabla^2 + k^2$}

La identidad de Green para el operador de Helmholtz da:

\begin{equation}
    \int_\Omega (\nabla^2 + k^2)\hat{\phi}\, d\Omega
    = \int_\Gamma \frac{\partial\hat{\phi}}{\partial n}\, d\Gamma
    + k^2 \int_\Omega \hat{\phi}\, d\Omega
\end{equation}

Nótese que \textbf{no se elimina} la integral de dominio a menos que $k = 0$,
ya que el termino $k^2\int_\Omega\hat{\phi}\,d\Omega$ persiste. Esto muestra
que el operador de Helmholtz \textbf{no es adecuado} para la conversion pura
de integrales de dominio que nos ocupa, aunque si lo es en el contexto de
resolver EDPs de Helmholtz donde este termino se absorbe en la formulacion
integral.

\cite{muleshkov1999helmholtz} derivan soluciones particulares para operadores
de tipo Helmholtz con splines poliharmonicos de orden superior, generalizando
los resultados de \cite{golberg1999annihilator} para thin plate splines.

\textbf{Compatibilidad con $b$:} optimo cuando $b$ tiene estructura oscilatoria
con numero de onda caracteristico $k$, o cuando se resuelve una EDP de Helmholtz
donde el termino fuente tiene esa naturaleza.

\subsection{Operador de Helmholtz Modificado: $\mathcal{L} = \nabla^2 - \lambda^2$}

La identidad es estructuralmente similar a la de Helmholtz:

\begin{equation}
    \int_\Omega (\nabla^2 - \lambda^2)\hat{\phi}\, d\Omega
    = \int_\Gamma \frac{\partial\hat{\phi}}{\partial n}\, d\Gamma
    - \lambda^2 \int_\Omega \hat{\phi}\, d\Omega
\end{equation}

El mismo problema persiste: hay un residuo de dominio. Las soluciones
particulares para este operador tienen interes en problemas de difusion-reaccion
donde $b$ tiene estructura de decaimiento exponencial con escala $\lambda^{-1}$.

\subsection{Operador Biharmonico: $\mathcal{L} = \nabla^4$}

La identidad de Green de cuarto orden es:

\begin{equation}
    \int_\Omega \hat{\phi}\,\nabla^4 u\, d\Omega
    = \int_\Gamma \left(
    \hat{\phi}\frac{\partial\nabla^2 u}{\partial n}
    - \frac{\partial\hat{\phi}}{\partial n}\nabla^2 u
    + \nabla^2\hat{\phi}\frac{\partial u}{\partial n}
    - u\frac{\partial\nabla^2\hat{\phi}}{\partial n}
    \right) d\Gamma
\end{equation}

Tomando $u = 1$ toda la integral de dominio se convierte en frontera, pero con
\textbf{cuatro terminos} en lugar de uno. La complejidad es mayor, pero la
conversion es exacta en frontera.

\cite{muleshkov2007multihelmholtz} desarrollan soluciones particulares para
operadores multi-Helmholtz de la forma
$\prod_{i=1}^m (\nabla^2 \pm k_i^2)$, que por el teorema de Hormander
\cite{tsai2009splines} cubren muchos sistemas de EDPs de ingenieria al
factorizar su operador determinante.

\textbf{Compatibilidad con $b$:} optimo cuando $b$ tiene estructura biharmonica,
por ejemplo en problemas de placas elasticas donde la deflexion satisface
$\nabla^4 w = q/D$.

\subsection{Tabla Comparativa}

\begin{center}
\begin{tabular}{llll}
\toprule
Operador $\mathcal{L}$ & Terminos en $\partial\Omega$ & Residuo dominio
    & Compatibilidad con $b$ \\
\midrule
$\nabla^2$ & 1 & No
    & $b$ suave, general \\
$\nabla^2 + k^2$ & 1 & Si ($k^2\int\hat\phi$)
    & $b$ oscilatoria, escala $k$ \\
$\nabla^2 - \lambda^2$ & 1 & Si ($\lambda^2\int\hat\phi$)
    & $b$ difusiva, escala $\lambda^{-1}$ \\
$\nabla^4$ & 4 & No
    & $b$ biharmonica \\
$\prod_i(\nabla^2 \pm k_i^2)$ & $2m$ & No (con $u=1$)
    & $b$ en espacio nulo del operador \\
\bottomrule
\end{tabular}
\end{center}

% ─────────────────────────────────────────────────────────────
\section{Hipotesis de Compatibilidad Optima}
% ─────────────────────────────────────────────────────────────

Con base en la teoria del aniquilador y la literatura revisada, pueden
formularse las siguientes hipotesis:

\begin{hypothesis}[Compatibilidad espectral]
\label{hyp:espectral}
Sea $\mathcal{A}$ el aniquilador de $b$, es decir $\mathcal{A}b = 0$. La
conversion de $\int_\Omega b\,d\Omega$ a integrales de frontera es
\textbf{exacta} cuando $\mathcal{L} = \mathcal{A}$, porque en ese caso la
base $\{f_j\} = \{\mathcal{L}\hat{\phi}_j\}$ puede reproducir exactamente
$b$ en el espacio nulo de $\mathcal{A}$.
\end{hypothesis}

\begin{hypothesis}[Minimalidad de terminos de frontera]
\label{hyp:minimalidad}
El operador optimo para la conversion pura de integrales de dominio (sin EDP
de fondo) es aquel de orden minimo cuya identidad de Green no genera residuos
de dominio al tomar $u = \text{const}$. El operador de Laplace $\nabla^2$ es
el unico operador de segundo orden que satisface esta condicion.
\end{hypothesis}

\begin{hypothesis}[Jerarquia poliharmonica]
\label{hyp:jerarquia}
Para $b$ con estructura que puede expresarse como combinacion de funciones en
el espacio nulo de un operador poliharmonico $\nabla^{2m}$ o de un producto
de operadores de Helmholtz $\prod_i(\nabla^2 \pm k_i^2)$, usar el operador
correspondiente como $\mathcal{L}$ minimiza el numero de nodos de colocacion
necesarios para una precision dada, porque las funciones base son exactamente
del tipo correcto para representar $b$.
\end{hypothesis}

\begin{remark}
Las hipotesis anteriores son consistentes con la literatura pero, hasta donde
es conocido, no existe una prueba formal de optimalidad que, dado solo $b(x)$
como funcion conocida y sin una EDP de fondo, permita seleccionar $\mathcal{L}$
de forma que minimice el error de cuadratura resultante. Esta es una pregunta
abierta de investigacion.
\end{remark}

% ─────────────────────────────────────────────────────────────
\section{El Enfoque de Hormander y los Productos de Helmholtz}
% ─────────────────────────────────────────────────────────────

Un resultado particularmente poderoso, discutido en \cite{tsai2009splines}, es
la aplicacion de la teoria de Hormander para operadores diferenciales lineales.
Muchos sistemas de EDPs de ingenieria tienen la propiedad de que el determinante
de su operador matricial es un producto de operadores de Helmholtz o
poliharmonicos:

\begin{equation}
    \det(\mathcal{L}_{ij}) = \prod_{p=1}^{m}(\nabla^2 \pm k_p^2)
\end{equation}

En estos casos, por el principio de Hormander, el sistema puede reducirse a
ecuaciones escalares con estos operadores, y las soluciones particulares para
los operadores de tipo Helmholtz y poliharmonico que se derivan en
\cite{muleshkov2007multihelmholtz} y \cite{tsai2009splines} cubren una clase
muy amplia de problemas de ingenieria.

El enfoque alternativo del \textbf{Radial Integration Method} (RIM) de
\cite{gao2002rim} es notable porque no usa $\mathcal{L}$ en absoluto: la
conversion a integrales de frontera se hace puramente por transformacion
geometrica en coordenadas esfericas, sin requerir que $b$ sea compatible con
ningun operador. Esto lo hace universalmente aplicable, aunque pierde la
eficiencia que se obtiene cuando hay compatibilidad entre $b$ y $\mathcal{L}$.

% ─────────────────────────────────────────────────────────────
\section{Consecuencias Practicas para el Curso}
% ─────────────────────────────────────────────────────────────

La discusion teorica anterior tiene consecuencias practicas concretas que
pueden resumirse en los siguientes principios de diseno:

\textbf{Principio 1 --- Por defecto, usar $\nabla^2$:} en ausencia de
informacion sobre la naturaleza de $b$, el operador de Laplace es la eleccion
mas robusta porque genera un solo termino de frontera, admite RBFs con
$\hat{\phi}$ en forma cerrada, y el sistema de colocacion es generalmente
bien condicionado.

\textbf{Principio 2 --- Explotar la estructura de $b$ si se conoce:} si $b$
es oscilatoria con escala caracteristica conocida, usar $\mathcal{L} = \nabla^2
+ k^2$ permite que las funciones base sean del tipo correcto. Si $b$ tiene
estructura biharmonica, usar $\nabla^4$.

\textbf{Principio 3 --- El aniquilador como guia:} si $b$ satisface una EDP
homogenea $\mathcal{A}b = 0$ (aunque sea aproximadamente), el operador
$\mathcal{A}$ es el candidato natural para $\mathcal{L}$.

\textbf{Principio 4 --- Verificar la eliminacion del dominio:} antes de usar
un operador dado, verificar que su identidad de Green con $u = \text{const}$
no genera terminos de dominio residuales. Solo $\nabla^2$ y $\nabla^{2m}$
(y sus productos con condiciones apropiadas) satisfacen esto.

\textbf{Pregunta de investigacion abierta:} dada unicamente $b(x)$ como
funcion numericamente conocida en un conjunto de puntos, ?`existe un criterio
a priori para seleccionar $\mathcal{L}$ que minimice el error de la cuadratura
resultante? La respuesta no es conocida en general.

\bibliographystyle{unsrt}
\bibliography{operador_b}

\end{document}
